\chapter{Conclusion}

\label{ch:conclusion}



This chapter summarises the findings discussed so far by tackling the research questions and outlining the directions for future research. 
Below are the answers to the two main research questions:

\subsection{How would the participants describe the experience of using voice input in the presented game environment?}


Throughout the study, each participant noted their own emotional responses to the voice input modality. 
Despite those differences, it can be observed that participants followed general patterns:
A large group of participants described the experience as deep, immersive and absorbing.
For some, the voice input allowed them to feel closer to someone in a role-playing session. For others, the open-ended nature of voice input created a feeling of higher control.

On the other hand, some participants reported feeling frustrated while playing the voice version.
Technical issues were the main reason for this feeling. Some players also felt a mismatch between their expectations and reality, adding to that feeling of frustration.


What stands out is the universal feeling of excitement. Each participant felt excited about using the voice input system. However, that emotion wasn't caused by the voice input itself, but rather the fact that it was a novel experience for all the participants.

\subsection{How do users feel about using the voice modality compared to traditional input?}

All of the participants felt a difference between the two versions of the game.
One of the biggest factors contributing to that was the differing amount of friction that they felt while playing the voice version.
Due to a feeling of reduced control, facing technical issues and VUI bugs, some participants felt frustrated and powerless. 
Despite the varying degrees of this feeling, all of the participants felt a higher amount of friction while playing the voice-driven game compared to the keyboard experience. 





Another big emotional difference was the feeling of immersion. 
Some of the players experienced a deeper trance-like state while playing the voice version. 
Using voice, needing to think more and sometimes formulating the commands in the first person all contributed to that feeling.
It should be noted that for another portion of players, immersion in the voice version was lower than in the keyboard version.
Players noted many factors, ranging from technical issues to a feeling of lack of control over the game, that contributed to that emotional state.
Again, all the participants felt a difference in immersion between the two versions; none remained neutral about it.


All of those findings indicate that participants felt differently between the two game versions. 
It also shows that comfort can be more important than immersion when it comes to game preference.



\section{Limitations}



It should be noted that the whole process was influenced by the Observer Effect. The sole fact that the participants knew that they were being observed had an impact on the results of the observations. This was especially evident when participants used their voice to navigate, since, despite any efforts from the author's side, the act of speaking to the computer feels more intimate compared to using a keyboard. This is not only  speculation but was directly mentioned by two of the participants, who implied that their emotions during playing would have been different had she not been under observation. 

Also, as discussed before, the novelty effect of being exposed to a new form of input had an emotional impact on participants.
They were all excited to try the voice input and alter their perception of it due to this being their first experience with voice-driven navigation in games.
This effect is hard to filter out completely, but the fact that participants had a clear preference split between the input versions suggests that it wasn't a factor that invalidated the whole study.





\section{Future research}

Future research about this topic could include trying to find the cause for the observed variance in the emotional impact of friction experienced in-game. The same goes towards the feelings of immersion - the lack of a clear model, the variance in those emotions is a question waiting to be answered.